\section*{A prescribed base residue and the descent from split short series}

Fix the base tuple $\mu _0$, its supplied record $O$ with the original
$\mathrm{Hk3.Core}$, tie and embedding pins, and
$K=O.c'.F$ with the specified embedding $\sigma:K\hookrightarrow
\mathbb C$. Write $n=\operatorname{rank}V_i$, $s=\operatorname{rank}W_i=D.R.r=r$, $b=O.\mathrm{bad}(())>0$,
$\mathcal O=\mathcal O_E$, and $\mathcal L=D_w.\mathrm{La}$.
We use the standard $V_i$ frame, the pinned $W_i$ lattice embedding
and $\mathrm{HalfIs}$ at this tuple. The finite datum in the spanning
argument has to do two jobs. Its
Fourier transform weights a geometric short series, and its untransformed
indicator has to be the precise finite oscillator datum on the base tuple.

% BEGIN S74B1
\paragraph{S74B1. A generating residue tuple with the prescribed finite indicator.}
Put $\mathfrak a=\prod_{P\in P_g}P$ and
$Q=\mathcal L/\bar{\mathfrak a}\mathcal L$.
The pinned half-level identifies $\bar{\mathfrak a}$ with
$M_{\rm fr}(l)$. There exists a tuple $x_1\in Q^n$ generating $Q$
over $\mathcal O$ for which the original finite test function at
$\mu _0$ is exactly $\mathbf1_{x_1}$ on residue tuples.

Indeed the base twisting character is trivial. Its ramification
criterion forces $Q(\mu _0)=\varnothing$, so its factor in the finite
weight is the empty product. Choose the supplied idempotent over
$P_g$. At $P\in P_g$, let $k_P=\mathcal O/P$ and let
$\rho_P:\mathcal L\to k_P^s$ be the reduction of the integral row
pairings in the $W_i$ frame. Its scalar rule is
$\rho_P(c x)=(\bar c\bmod P)\rho_P(x)$. We verify that the joint
map onto $\prod_P k_P^s$ has kernel $\bar{\mathfrak a}\mathcal L$.

Here are the denominator details for the lattice embedding $\iota$.
Let $L$ denote the self-dual frame lattice, $N=\operatorname{Nm}(\mathrm{ccK})$
and $m=\operatorname{Nm}(\mathfrak b)$ in the host data.
The defining lattice inclusions give
$N L\subset\iota(\mathcal L)$ and $m\iota(\mathcal L)\subset L$.
The half-level primes avoid $Nm$ by the coprimality with the pinned
bad factors and the half-different condition. Use finite dual
generators $s'_j$ of $L$. Their row pairing matrix $B$ has an integral
adelic left inverse $A$ with $AB=1_s$. More explicitly a finite
projective dual family resolves the identity on $L$; in frame
coordinates its evaluation rows give $B$ and its coefficient columns
give $A$. The coefficients $c_j$ of $y=\sum_j\bar c_j s'_j$ are the
entries of its pairing row times $A$. Choose $z_j\in\mathcal L$ with
$\iota z_j=Ns'_j$. For prescribed rows $t_P$, lift the conjugate
coefficients of a linear combination of the $z_j$ by the Chinese
remainder theorem, prescribing the row $N^{-1}t_PA_P$ at $P$.
The resulting pairing rows are $t_PA_PB_P=t_P$.

If all pairings of $x$ vanish modulo the indicated primes, the
coefficients $c_j$ of $m\iota x$ vanish at every $P$ by the same
coefficient formula. Thus $Nm\,x=\sum_j\bar c_j z_j$ is in
$\bar{\mathfrak a}\mathcal L$. Pick $v\in\mathcal O$ with
$Nm v\equiv1\pmod{\mathfrak a}$. The identity
$x=\overline{1-Nm v}\,x+\bar v(Nm\,x)$ gives the desired inclusion.
The reverse inclusion is the scalar rule. This proves the kernel and
surjectivity statements including the case of an empty set of primes.

Let $U_P$ be the reduced $V_i$ frame matrix. Both it and its inverse
are integral before reduction. Since $s\le n$, write
$J=\binom{1_s}{0}$. Lift the rows of $U_PJ$ simultaneously through
the joint map, and call their classes in $Q$ the entries of $x_1$.
For any pairing row $t_P$, the row $(t_P,0)U_P^{-1}$ gives
coefficients expressing $t_P$ in the lifted rows at $P$.
Conjugate Chinese remainder lifts make these coefficients simultaneous,
so the entries span $Q$. For an arbitrary tuple with matrix of rows
$T_P$, its residue equals $x_1$ exactly when $U_P^{-1}T_P=J$ at
every $P$. These are the upper identity and lower zero tests defining
the finite oscillator factor at the base. Both its values, zero and
one, are preserved under $\sigma$.
% END S74B1

% BEGIN S74B2
\paragraph{S74B2. A field of definition for the short-series move.}
Suppose a rational unitary move carries $(\Lambda,N,P,w)$ to
$(\Lambda',N',P',w')$, with the integral harmonic, level and
connection hypotheses of the original short-value construction.
Take $P'$ to be the pushforward of $P$ and the weights to agree on
moved residues. Write $U,U'$ for the two normalization-one short
values. They are natural functions of framed points by T73C1.
The analytic change-of-vacuum input on the pinned complex charts
takes the following precise form, with a single $|\tau|=1$:
\[
 U'(y)\sqrt{a}=\tau\sqrt{a'}\, U(y),\qquad
 a=[\Lambda:\Lambda\cap\mathrm{Std}]^{p+q},\quad
 a'=[\Lambda':\Lambda'\cap\mathrm{Std}]^{p+q}.
\]
It suffices to give this identity on analytic representatives for frame
changes since $P,P'$ have the same weight. Then there is one $z\in K$
such that for \emph{every} commutative $K$-algebra $S$ and framed
$S$-point $y$, $U(y)=zU'(y)$ (with $z$ mapped into $S$).

To see this, the indices are positive and hence the complex ratio is
$c=\sqrt{a}/(\tau\sqrt{a'})$. If $U'$ vanishes at all complex
points, both readings vanish there and we take $z=1$. Otherwise
choose a point where $U'$ is nonzero. For each $K$-automorphism
$\alpha$ of $\mathbb C$, transport the frame by $\alpha$. Naturality
and the identity at its transported point give
$(\alpha(c)-c)\alpha(U'(y))=0$. Thus $c\in\sigma(K)$. Here the
fixed field is $\sigma(K)$: a separable algebraic element outside it
has a different conjugate, and a transcendental element can be
included in a transcendence basis and translated by one. The resulting
isomorphisms extend to the algebraic closure. Take $z=\sigma^{-1}(c)$.

For the extension to test algebras, use the framed chart cover of
TDR3 and TDR5--TDR7; D74U2 supplies its section law at power one.
Over each maximal localization $S_{\mathfrak m}$ a chart containing
the image of the closed point contains the entire image. The universal
frame of its reduced finite type $K$-chart $A$ and the given frame
therefore pull back to a common frame via $A\to S_{\mathfrak m}$
and $S\to S_{\mathfrak m}$.
On $A$ the difference $U-zU'$ dies under every $A\to\mathbb C$.
Each maximal residue field is finite over $K$ and embeds in
$\mathbb C$ over $K$. The intersection of the maximal ideals is the
nilradical, zero in $A$. The difference is therefore zero in $A$.
Naturality gives equality in every maximal localization of $S$, hence
in $S$ itself. This explains why the argument permits nonreduced $S$.
The analytic vacuum identity is the remaining input to this descent.
% END S74B2

% BEGIN S74B3
\paragraph{S74B3. From the identity cover to a single generating residue.}
Assume the stable short-series span on the pinned identity cover, at
$n\ge 4[E:\mathbb Q]r+4$, and the complex vacuum identity in S74B2 for
rational unitary transport. Every cusp reading then lies in the complex span of raw
sections with the \emph{one} weight $w=\widehat{\mathbf1}_{x_1}$ of
S74B1. Each such section uses a unimodular $\Lambda$ trivial at the
bad modulus and at $l.M$, a permitted $N$, and
$P\in\mathrm{harmL}(K,\Lambda)$ of the given weight.

In fact an identity-cover generator with lifted tuple $\tilde x_0$
has normalized short value $f$ with scalar $\beta_H(\Lambda)\ne0$.
Its multiplier is one on the identity cover. Evaluating there on
geometric fibers and using uniqueness shows that $f$ has the required
normalized value on every reduced test chart (geometric points
separate elements of that chart). Multiplying the defining short
identity by $\beta_H$ gives the raw normalization-one value for
$\beta_H f$. The Fourier sum against the point indicator reduces to
its single summand, the level character of $\tilde x_0$; this
character depends only on its residue $x_0$. Admissibility of the
lifted tuple says precisely that its entries generate after passage
to the quotient. Hence allowing these $x_0$ already spans all the
cusp readings in question. The numerical bound in the supplied
record is exactly this bound since $[E:\mathbb Q]=2f$.

To construct the move, work first on the residue module of S74B1.
On each factor both tuples give surjections from an $n$-dimensional
space to the same $s$-dimensional space. Choose bases of their
kernels and lifts of a common basis of the target. The change of
these bases carries one surjection to the other. In a split pair the
unitary matrix is specified by its matrix on one factor, with inverse
transpose on the conjugate factor. Prescribe the first factor so the
adjoint sends the tuple $x_1$ to $x_0$; the two prescriptions are
compatible by construction. Lift the invertible residue matrices to
integral invertible local matrices and choose identity matrices at
the other places over $2b$. Weak approximation in the rational
unitary group gives a rational unitary matrix in these finitely many
open congruence neighborhoods (use the rational Cayley chart after
perturbing the local matrices away from its excluded divisor).
Its action is integral with integral inverse at these places. For
transported residues $q,q'$ the level pairings satisfy
$\langle x_1,q'\rangle=\langle x_0,q\rangle\pmod M$ by the
adjoint identity. Exponentiating proves precisely equality of the
Fourier weights. The resulting matrix preserves the integral
harmonic module by pushforward, and
triviality at $2$ and at the level supplies an allowable $N'$.
If $s_0$ is the old raw section, S74B2 gives $s_0=zU'$ on every test
chart. For $z=0$ its reading is zero. Otherwise $z^{-1}s_0$ is still
a section of the same weight, and the displayed identity says at
every test point that it is a raw section for the new tuple. Thus its
$K$-span contains $s_0$. This does not require constructing a
section from an arbitrary function on complex points. T73C2--T73C3
supply the original pinned short rules used on the identity cover.
The stable analytic span used here still needs the reproducing
kernel and level-to-short-series comparisons specified in ICL14 and
ICL16.
% END S74B3

% BEGIN S74B4
\paragraph{S74B4. Moving a raw lattice to a split finite point without moving its weight.}
Let $b$ be the supplied bad integer, $M=l.M$, $M\mid b$, and
$2\mid b$. If $\Lambda$ is any of the raw lattices in S74B3,
there are a pure finite $h$ supported on split places away from $bM$
and a rational unitary $d$ such that
\[
       d\Lambda=\overline{L_h},\qquad
       L_h=\mathrm{hLat}(V_i,h),\qquad d\equiv1\pmod{bM}
\]
in the columnwise congruence sense (denominators coprime to $bM$).
In particular the Fourier weight $\widehat{\mathbf1}_{x_1}$ can be
used on both sides of this move.

For clarity construct first a finite adelic unitary $g$ carrying the
standard completed lattice to $\widehat{\bar\Lambda}$ and equal to
one at primes over $bM$. Conjugation preserves the standard lattice,
unimodularity and level triviality. The two lattices coincide at the
places over $bM$ by triviality and at all but finitely many places by
commensurability. At the remaining places use an integral orthonormal
frame of the self-dual lattice. Here the host gives an unramified CM
extension at finite places. At a split place use a basis in one factor
and its Hermitian dual in the other. At an unramified quadratic field
place the residual nondegenerate Hermitian form contains a vector of
nonzero norm: if all norms vanished, polarization with all residue
scalars and nondegeneracy of the residue trace pairing would kill the
entire form. Lift this vector. Its norm is a unit, and the norm on
units is surjective by the finite-field norm followed by successive
correction on principal units using the residue trace. Rescale to norm
one, split off the orthogonal line, and induct. This gives the required
local frames and hence $g$.

Apply the finite factorization (BTH.49) at modulus $bM$, writing
$g=q_f h k$ with $q$ rational unitary and $k$ in the standard
lattice stabilizer congruent to one at the modulus. The split support
of $h$ gives $h=1$ there. Consequently $q\equiv1\pmod{bM}$
there, hence also as a rational matrix after clearing denominators
coprime to $bM$. Since $k$ preserves the completed standard lattice,
intersecting the lattice identity with the rational space gives
$\bar\Lambda=qL_h$. Unitarity gives $\bar q^{-1}=q^{\mathsf T}$,
so $d=q^{\mathsf T}$ has precisely the claimed direction.
The residue action of $d$ modulo $M$ is the identity, so any function
of the residue, including this Fourier weight, is preserved.

The split-support lattice argument of RTG7 supplies a permitted $N'$
and the original lattice guard for $L_h$; the conjugation in the last
display agrees with the lattice in the short value. The change of
variables pushes $P$ to its integral harmonic polynomial on
$d\Lambda$: pullback on linear letters commutes with contractions by
the Hermitian form and identifies the two coefficient lattices.
Under the analytic change-of-vacuum input of S74B2, the resulting
short value $t$ on the original dictionary thus satisfies
$L s_0=\sigma(z)t$ for the original raw section $s_0$,
where $L$ is its reading followed by evaluation on that dictionary.
% END S74B4

% BEGIN S74B5
\paragraph{S74B5. Reduction of the base class-sum span to the original analytic comparisons.}
Keep the full pinned base record. Let $\mathcal S$ be its cusp
sections for which some coefficient tuple satisfies both the original
place condition and the original geometric class-sum rule. Assume
the stable span and transport inputs of S74B3 and the two point
comparisons below. Then every cusp section lies in
$\operatorname{span}_K\mathcal S$, as required by
$\mathrm{ClassSumsSpan}$.

Write $\ell=\mathrm{lieFac}>0$ and
$a(g)=c_{\mathrm{dict}(g)}\ell\,\mu(g)$, using the original frame
constants and unitary character. For a moved lattice $\bar L_h$,
let $t(g)$ be its normalization-one short value. The harmonic pin
expresses its polynomial as $O.\mathrm{io}(v)$ with $v$ in the
prescribed coefficient space. The needed one-class kernel comparison
is
\[
     K_{h,v}(g)=A_h\,a(g)t(g),\qquad
 A_h=\mathrm{cmpConst}(D,k_a,\delta,\deg W)\,
         \ell^{-1}\sigma(\mathrm{epF}(L_h)).
\]
The constants can be divided out: $k_a$ has norm one, $\delta$ is a
sign, and the lattice guard and the split determinant ideal
$\mathfrak c$ give
$\mathrm{epF}(L_h)=B(\mathfrak c)^{-1}q$ with $q>0$ rational.
The ideal is nonzero and coprime to the bad modulus; the pinned value
of $\sigma(B(\mathfrak c))$ is a unitary value times a positive
square root of its norm. In addition $a(g)\ne0$ for each $g$.

By K73G3, $K_{h,v}=\sum_j b_j\Theta_j$ for finitely many lifts of
actual rational form inputs. For each input choose its base-class
coordinates by the pinned representative presentation. This witnesses
the original place condition. The geometric class-sum construction of T73C5
furnishes a section $f_j\in\mathcal S$ on these coordinates.
The second needed point comparison, in the \emph{same} frame
normalization as the first, is $\Theta_j(g)=u\,a(g)Lf_j(g)$,
where $u$ is the common norm-one phase for the base record. Substitution
and pointwise cancellation of $a(g)$ give
\[
       t=\frac{u}{A_h}\sum_j b_j Lf_j .
\]
This calculation allows zeros of $t$ and zeros of individual lifts.
S74B4 places every raw reading in the same span, and S74B3 places
every cusp reading there. Finally K73G4 descends this finite span
through $L$ to the coefficient field $K$ at the original cusp
sections.

This is a reduction of the base spanning requirement. The stable
reproducing kernel with its chart comparisons, the analytic
change-of-vacuum identity, and the geometric comparisons on the same
pinned normalization are substantive inputs. Finite-point generation
and coefficient descent alone give no identity for those inputs.
% END S74B5
